% Part: methods
% Chapter: proofs
% Section: using-definitions

\documentclass[../../../include/open-logic-section]{subfiles}

\begin{document}

\olfileid{mod}{prf}{def}

\olsection{د تعريفونو کارول}

مخکې مو وويل چې له هغو ټولو تعريفونو سره بلد اوسئ چې
په ثبوت کښې کارېدلے شي، او سم يې وکارولے شئ۔ دا ډېر
مهم ټکے دے؛ ښه ده چې لږ نور يې هم وڅېړو۔ تعريفونه
ځانګړنې او اړيکې لنډوي، څو پرې په لنډه توګه خبرې
وکړو۔ هغه نوے راوړل شوے لنډ نوم \emph{تعريفېدونکے}
(\texttt{definiendum}) بلل کېږي، او هغه څه چې نوم يې
لنډوي \emph{تعريفوونکے بيان} (\texttt{definiens}) دے۔
په ثبوتونو کښې ډېری وخت
بايد بېرته وګورو چې تعريفېدونکے څنګه راوړل شوے و،
ځکه د ثبوت د پرمخ بيولو دپاره د تعريفوونکي بيان
منطقي جوړښت کاروو: هماغه اوږده بڼه چې ټاکلے نوم
يې لنډيز دے۔ د تعريفونو په پرانيستلو سره د استدلال
اصل منطقي جوړښت څرګندوئ۔

راځئ په يوه بېلګه پيل وکړو۔ فرض کړئ دا خبره ثابتول غواړئ:

\begin{prop}
د هر دوو سټونو $A$ او $B$ دپاره، $A \cup B = B \cup A$۔
\end{prop}

د ثبوت له پيلولو مخکې بايد پوه شو چې د دوو سټونو
يو شان کېدل څه معنا لري؛ يعنې په دې مساوات کښې د
سټونو دپاره «$=$» څه معنا لري۔ سټونه هغه وخت يو
شان تعريفېږي چې يو شان !!{element}s ولري۔ نو دا
تعريف بايد پرانيزو:

\begin{defn}
سټونه $A$ او $B$ \emph{يو شان} دي، $A = B$، که او يوازې
که د~$A$ هر !!{element} د~$B$ !!a{element} وي، او برعکس۔
\end{defn}

دا تعريف $A$ او~$B$ د هر ډول سټونو د ځايناستو
نښو په توګه کاروي۔ هغه څه چې تعريفوي --- \emph{تعريفېدونکے}
--- «$A = B$» عبارت دے؛ تعريف هغه شرط ښيي چې
له مخې يې $A = B$ رښتيا وي۔ همدا شرط --- «د~$A$ هر
!!{element} د~$B$ !!a{element} وي، او برعکس» ---
\emph{تعريفوونکے بيان} دے۔\footnote{په همدې ځانګړې
  بېلګه کښې --- که څه هم دا خبره مغشوشوونکې ښکاري! ---
  کله چې $A = B$ وي، سټونه $A$ او $B$ په رښتيا يو
  او هماغه سټ وي، سره له دې چې د مساوات په کيڼ
  او ښي اړخ کښې يې په بېلو تورو ښيو۔ خو د هماغه
  سټ د پېژندلو لارې بېلې کېدے شي؛ له همدې امله
  تعريف بې ګټې نه دے۔} تعريف ټاکي چې $A = B$
هغه وخت، او يوازې هغه وخت، رښتيا دے چې ياد شرط
پوره شي؛ «که او يوازې که» ته په لنډه \texttt{iff} هم وايي۔

د تعريف د کارولو پر مهال بايد د تعريف $A$ او $B$
له خپلې ځانګړې بېلګې سره برابر کړئ۔ زموږ په بېلګه
کښې، د $A \cup B = B \cup A$ د رښتيا کېدو دپاره د
$z \in A \cup B$ هر غړے بايد په $B \cup A$ کښې هم
وي، او برعکس۔ په بيان کښې $A \cup B$ د تعريف
د~$A$ رول لري، او $B \cup A$ د~$B$ رول۔ ځکه چې
$A$ او $B$ د تعريف او د ثابتېدونکي بيان په دواړو
ځايونو کښې راغلي، خو بېل رولونه لري، پام وکړئ
چې سره يې ګډ نه کړئ۔ د بېلګې په توګه، دا به
تېروتنه وي چې بيان د دې په ښودلو ثابت کړئ چې
د~$A$ هر !!{element} د~$B$ !!a{element} دے، او
برعکس: دا به $A = B$ وښيي، نه دا چې
$A \cup B = B \cup A$۔ (څرنګه چې $A$ او $B$
هر دوه سټونه کېدے شي، له $A$ او $B$ په اړه
بې له نورو فرضونو هغوی په رښتيا بېل هم
کېدے شي؛ نو دغه لار
مو وړاندې نه بيايي۔)

په ثبوت کښې د سټونو د نظريې مفاهيم، لکه اتحاد،
کاروو؛ نو د ثبوت د لارې د پوهېدو دپاره د $\cup$
نښې په معنا هم پوهېدل پکار دي۔ کله کله د يوه
تعريف له پرانيستلو وروسته نور تعريفونه هم
پرانيستل غواړي۔ د بېلګې په توګه، $A \cup B$
د $\Setabs{z}{z \in A
  \text{ or } z \in B}$ په توګه تعريف شوے دے۔ نو که
غواړئ $x \in A \cup B$ ثابت کړئ، د $\cup$ تعريف
پرانيستل درته وايي چې
$x \in \Setabs{z}{z \in A \text{ or } z \in B}$
ثابت کړئ۔ اوس دا هم په ياد ولرئ چې
$x \in \Setabs{z}{\dots z\dots}$ هغه وخت او يوازې
هغه وخت وي چې $\dots
x\dots$۔ له همدې امله، د
$\Setabs{z}{\dots z \dots}$ نښې تعريف نور هم
پرانيزو: بايد وښيو چې $x \in A$ يا $x \in B$۔
نو «د $A \cup B$ هر
!!{element} د $B \cup A$ !!a{element} هم دے»
په رښتيا دا معنا لري: «د هر $x$ دپاره، که $x \in
A$ يا $x \in B$ وي، نو $x \in B$ يا $x \in A$
وي»۔ که د بيان ټول تعريفونه پرانيزو، دا
راڅرګندېږي چې لاندې خبره بايد ثابته کړو:

\begin{prop}
د هر دوو سټونو $A$ او $B$ دپاره: (الف) د هر $x$
دپاره، که $x \in A$ يا $x \in B$ وي، نو $x \in B$
يا $x \in A$ وي؛ او (ب) د هر $x$ دپاره، که
$x \in B$ يا $x \in A$ وي، نو $x \in A$ يا
$x \in B$ وي۔
\end{prop}

مهمه خبره دا ده چې د تعريفونو پرانيستل د ثبوت
د جوړولو اړينه برخه ده۔ سم کارول يې کله کله
ګران وي: د تعريف متغيرونه او د ثابتېدونکې دعوې
عبارتونه بايد سره بېل کړئ او په دقت يې سره
برابر کړئ۔ د برياليتوب دپاره بايد پوه شئ چې
پوښتنه څه غواړي او په هغې کښې هر کارېدلے
عبارت څه معنا لري؛ ډېر ځله به تر يوه زيات
تعريف پرانيستل غواړئ۔ په لاندې ساده ثبوتونو
کښې حل نږدې په نېغه له تعريفونو راځي۔ البته
تل به خبره دومره ساده نه وي۔

\begin{prob}
فرض کړئ له تاسو څخه غوښتل کېږي چې
$A \cap B \neq \emptyset$ ثابت کړئ۔ دلته راغلي
ټول تعريفونه پرانيزئ؛ يعنې خبره داسې بيا
ووايئ چې «$\cap$»، «=»، يا «$\emptyset$»
پکښې ياد نه شي۔
\end{prob}

\end{document}
